\chapter{Geometric Brownian motion}
\label{ch:gbm}

\begin{question}
Stock prices cannot go negative and move by percentages, not by fixed amounts. Which process models them, and what is its expected value?
\end{question}

\begin{whatfor}
\begin{itemize}
  \item Defines geometric Brownian motion $X(t)=e^{Y(t)}$, with $Y$ a Brownian motion with drift.
  \item Computes $\E_P\big(X(t)\given\F_s\big)$, the mean, the variance, and when it is a martingale.
  \item Explains why i.i.d.\ percentage returns lead to it.
\end{itemize}
\skipline{fs:gbm}
\end{whatfor}

\begin{prereq}
From Part~0: $\E_Pe^{\xi}=e^{\mu+\sigma^2/2}$ for normal $\xi$, CE2--CE3. From \cref{ch:first-mg,ch:bm-walks,ch:drift}: Wiener process, independent increments, Brownian motion with drift and with variance parameter, central limit theorem, geometric random walk.
\end{prereq}

%=====================================================================
\section{Definition and conditional mean}

\begin{definition}[Geometric Brownian motion]
\label{def:gbm}
Let $Y(t)=\mu t+\sigma W(t)$ be a Brownian motion with drift $\mu$ and variance parameter $\sigma^2$, so $Y(t)\sim N(\mu t,\sigma^2t)$. The process
\[
  X(t)=X(0)\,e^{Y(t)}
\]
is a \dfn{geometric Brownian motion} (GBM). It is always positive, and $\ln X(t)-\ln X(0)=Y(t)$ is normal: $X(t)$ is \dfn{log-normal}.
\end{definition}

\begin{theorem}[Conditional mean of GBM]
\label{thm:gbm-mean}
With $\F_s=\sigma\big(X(u),0\le u\le s\big)$ and $s<t$,
\[
  \E_P\big(X(t)\given\F_s\big)=X(s)\,e^{(t-s)(\mu+\frac12\sigma^2)},\qquad
  \E_P\big(X(t)\big)=X(0)\,e^{t(\mu+\frac12\sigma^2)} .
\]
\end{theorem}
\begin{proof}
Split off the future increment of $Y$, as always.
\begin{steps}
  \item $X(t)=X(s)\,e^{Y(t)-Y(s)}$. Knowing $X(u)$ for $u\le s$ is knowing $Y(u)$ for $u\le s$, and $Y(t)-Y(s)$ is independent of that.
  \item CE2, CE3: $\E_P\big(X(t)\given\F_s\big)=X(s)\,\E_P\,e^{Y(t)-Y(s)}$.
  \item $Y(t)-Y(s)\sim N\big(\mu(t-s),\sigma^2(t-s)\big)$, so by the exponential moment its mean is $e^{\mu(t-s)+\frac12\sigma^2(t-s)}$.
  \item Take $\E_P$ with $s=0$ for the unconditional mean. \qedhere
\end{steps}
\end{proof}

\begin{result}[Variance; martingale condition]
\label{res:gbm-var}
With $X(0)=x_0$:
\begin{enumerate}[label=(\alph*)]
  \item $\Var X(t)=x_0^2\,e^{2\mu t+\sigma^2t}\big(e^{\sigma^2t}-1\big)$;
  \item $X$ is a martingale if and only if $\mu=-\frac12\sigma^2$; it is a submartingale if $\mu>-\frac12\sigma^2$ and a supermartingale if $\mu<-\frac12\sigma^2$.
\end{enumerate}
\end{result}
\begin{proof}
(a) $X(t)^2=x_0^2e^{2Y(t)}$ with $2Y(t)\sim N(2\mu t,\,4\sigma^2t)$. So $\E_PX(t)^2=x_0^2e^{2\mu t+2\sigma^2t}$; subtract $(\E_PX(t))^2=x_0^2e^{2\mu t+\sigma^2t}$. (b) The factor $e^{(t-s)(\mu+\sigma^2/2)}$ in \cref{thm:gbm-mean} is $1$, $>1$ or $<1$ accordingly.
\end{proof}

\begin{intuition}
Even if the exponent $Y$ has no drift, $X=e^Y$ drifts upward on average: up-moves of $e^Y$ are bigger than down-moves ($e^{1}-1>1-e^{-1}$). The correction $\mu=-\frac12\sigma^2$ cancels exactly this inflation; with it, $X(t)=x_0e^{\sigma W(t)-\frac12\sigma^2t}$ is the exponential martingale of \cref{ch:first-mg}.
\end{intuition}

\begin{pic}
\begin{tikzpicture}
\begin{axis}[width=10cm,height=4.8cm,xmin=0,xmax=1,ymin=0,xlabel={$t$},ylabel={$X(t)$},legend style={font=\scriptsize,at={(0.02,0.98)},anchor=north west},legend columns=2]
  \addplot[defc,thick] table[x=t,y expr={exp(0.5*\thisrow{t}+0.6*\thisrow{w1})}] {figs/bm3.dat}; \addlegendentry{paths}
  \addplot[defc,thick,forget plot] table[x=t,y expr={exp(0.5*\thisrow{t}+0.6*\thisrow{w2})}] {figs/bm3.dat};
  \addplot[defc,thick,forget plot] table[x=t,y expr={exp(0.5*\thisrow{t}+0.6*\thisrow{w3})}] {figs/bm3.dat};
  \addplot[trapc,very thick,dashed,domain=0:1] {exp(0.68*x)}; \addlegendentry{$\E_PX(t)=e^{0.68t}$}
\end{axis}
\end{tikzpicture}
\picnote{Three GBM paths with $x_0=1$, $\mu=0.5$, $\sigma=0.6$. They stay positive; the mean grows at rate $\mu+\frac12\sigma^2=0.68$, faster than the drift $\mu=0.5$ of the exponent.}
\end{pic}

%=====================================================================
\section{Why stock prices are modelled this way}

Let $X_n$ be a stock's price on day $n$. Suppose the daily ratios $Y_n=\frac{X_n}{X_{n-1}}$ are i.i.d.: today's percentage change does not depend on the past and has the same law every day. Then
\[
  X_n=Y_nY_{n-1}\cdots Y_1\,X_0,\qquad \ln X_n=\sum_{i=1}^n\ln Y_i+\ln X_0 .
\]
The $\ln Y_i$ are i.i.d., so $\ln X_n$ is a random walk; suitably scaled (central limit theorem, \cref{ch:bm-walks}) it approximates a Brownian motion with drift. Hence $X_n$ approximates a geometric Brownian motion. This is the discrete version: the geometric random walk of \cref{ch:first-mg}.

\begin{example}[An expected price]
\label{ex:gbm-price}
$X(t)=5e^{Y(t)}$ with $\mu=0.01$ and $\sigma=0.3$. Then
\[
  \E_P\,X(10)=5\,e^{10(0.01+0.045)}=5e^{0.55}\approx8.67,
\]
while the median price is $5e^{\mu\cdot10}=5e^{0.1}\approx5.53$: half the paths end below it.
\end{example}

\begin{remark}[Two meanings of ``drift'']
\label{rem:two-mu}
In Chapter~16, GBM reappears as the solution of $dX_t=\mu X_t\,dt+\sigma X_t\,dW_t$. \emph{There} $\mu$ is the drift of $X$ itself and $\E_PX(t)=x_0e^{\mu t}$. \emph{Here} $\mu$ is the drift of the exponent $\ln X$. The two are linked by $\mu_{\text{here}}=\mu_{\text{there}}-\frac12\sigma^2$. Always check which $\mu$ a question means.
\end{remark}

%=====================================================================
\formulasheet{fs:gbm}

\begin{mastercard}[label=card:gbm]{Geometric Brownian motion: everything from one normal}
$Y_t=Y_0e^{mt+\sigma W_t}$. Every question about $Y_t$ is a question about the normal variable $\ln Y_t$; answer it with \cref{card:normal-mgf} (moments) or \cref{card:normal-prob} (probabilities):
\[
  \boxed{\;\ln Y_t\given\F_s\ \sim\ N\big(\ln Y_s+m(t-s),\ \sigma^2(t-s)\big)\qquad(s=0:\ N(\ln Y_0+mt,\ \sigma^2t)).\;}
\]
\textbf{Which $m$?} Told ``$Y_t=Y_0e^{\mu t+\sigma W_t}$'' or ``the log of the price is BM with drift $\mu$'': $m=\mu$. Told the SDE $dY=\mu Y\,dt+\sigma Y\,dW$ or ``expected return $\mu$'': $m=\mu-\frac12\sigma^2$.
\begin{plugrows}
  \plugrow{Mean}{\pstep{Box with $s=0$: $\ln Y_t\sim N(\ln Y_0+mt,\sigma^2t)$.}\pstep{$Y_t=e^{\ln Y_t}$: \cref{card:normal-mgf} with $u=1$: $\E_Pe^{\xi}=e^{\mu_\xi+\sigma_\xi^2/2}$.}\pstep{Substitute $\mu_\xi=\ln Y_0+mt$, $\sigma_\xi^2=\sigma^2t$.}}{$\E_PY_t=Y_0e^{(m+\sigma^2/2)t}$}{read $m$ correctly}{$Y_0=5$, $\sigma=0.3$, $t=10$; exponent form $m=0.01$: $5e^{(0.01+0.045)10}=5e^{0.55}\approx8.67$; SDE form $\mu=0.01$, $m=0.01-0.045$: $5e^{0.01\cdot10}=5e^{0.1}\approx5.53$}
  \plugrow{Variance}{\pstep{\cref{card:Eh}, row 3: $\Var Y_t=\E_PY_t^2-(\E_PY_t)^2$.}\pstep{$Y_t^2=e^{2\ln Y_t}$: \cref{card:normal-mgf} with $u=2$: $Y_0^2e^{2mt+2\sigma^2t}$.}\pstep{Row 1 squared: $Y_0^2e^{2mt+\sigma^2t}$.}\pstep{Subtract and factor.}}{$\Var Y_t=(\E_PY_t)^2\big(e^{\sigma^2t}-1\big)$}{either form, once $m$ is read}{$Y_0=1$, $m=0$, $\sigma=1$, $t=1$: $\E_PY_1=e^{0.5}$, $\Var=e^{1}(e^1-1)\approx2.718\cdot1.718\approx4.67$}
  \plugrow{Conditional moments}{\pstep{Box: $\ln Y_t\given\F_s\sim N\big(\ln Y_s+m(t-s),\sigma^2(t-s)\big)$.}\pstep{$Y_t^n=e^{n\ln Y_t}$: \cref{card:normal-mgf} with $u=n$.}\pstep{Substitute: $e^{n\ln Y_s+nm(t-s)+\frac12n^2\sigma^2(t-s)}$.}}{$\E_P(Y_t^n\given\F_s)=Y_s^n\,e^{(nm+\frac12n^2\sigma^2)(t-s)}$; $n=1$: $Y_s\,e^{(m+\sigma^2/2)(t-s)}$}{$s<t$}{$Y_0=100$, $m=0.01$, $\sigma=0.2$, $n=1$: $\E_P(Y_{10}\given\F_0)=100e^{(0.01+0.02)10}=100e^{0.3}\approx134.99$}
  \plugrow{Martingale?}{\pstep{\cref{card:onestep}: martingale $\iff\E_P(Y_t\given\F_s)=Y_s$.}\pstep{Row 3, $n=1$: $\E_P(Y_t\given\F_s)=Y_se^{(m+\sigma^2/2)(t-s)}$.}\pstep{Equal to $Y_s$ $\iff$ the exponent is $0$.}}{martingale $\iff m+\frac12\sigma^2=0$; sub if $>0$; super if $<0$}{$s<t$}{$\sigma=0.2$: $m=-\frac12(0.04)=-0.02$; SDE form $m=\mu-\frac12\sigma^2$: $m+\frac12\sigma^2=\mu=0$}
  \plugrow{Median}{\pstep{Box with $s=0$: $\ln Y_t$ is normal, symmetric about $\ln Y_0+mt$.}\pstep{So $P(\ln Y_t\le\ln Y_0+mt)=\frac12$; exponentiate.}\pstep{Compare with row 1: the mean has the extra factor $e^{\sigma^2t/2}>1$.}}{median $Y_0e^{mt}$, below the mean}{$\sigma>0$}{$Y_0=100$, $m=0.08$, $\sigma=0.2$, $t=1$: median $100e^{0.08}=108.33$; mean $100e^{0.08+0.02}=110.52$}
  \plugrow{Probability above a level}{\pstep{Box with $s=0$: $\ln Y_t\sim N(\ln Y_0+mt,\sigma^2t)$.}\pstep{$\{Y_t>K\}=\{\ln Y_t>\ln K\}$.}\pstep{\cref{card:normal-prob}, row 2: $\Phi\big(\frac{\text{mean}-\ln K}{\text{s.d.}}\big)$.}}{$P(Y_t>K)=\Phi\Big(\dfrac{\ln(Y_0/K)+mt}{\sigma\sqrt t}\Big)$}{$K>0$}{same, $K=110$: $\frac{\ln(100/110)+0.08}{0.2}=\frac{-0.0953+0.08}{0.2}=-0.0766$; $\Phi(-0.0766)\approx0.470$}
  \plugrow{Given known values of the price}{\pstep{$\ln Y_t=\ln Y_0+mt+\sigma W_t$ is the process of \cref{card:Ylaw} with $x_0=\ln Y_0$, $\mu=m$.}\pstep{Apply \cref{card:Ylaw} to $\ln Y$ with the known $\ln Y_s$ values.}\pstep{Answer every question about $\ln Y$, then exponentiate.}}{e.g.\ $\ln Y_t\given\ln Y_s\sim N\big(\ln Y_s+m(t-s),\sigma^2(t-s)\big)$}{work with $\ln Y$, never with $Y$}{$m=0$, $\sigma=0.2$, $Y_1=110$: $\ln Y_2\given Y_1\sim N(\ln110+0,\,0.04)$; $P(Y_2>110)=P(\ln Y_2>\ln110)=\frac12$}
\end{plugrows}
\end{mastercard}
